\documentclass{handout}

\assignment{Homework 6}
\title{Energy: a new perspective}

\begin{document}
\maketitle
\practiceinstructions

\section{Work and the work--energy theorem}

\begin{questions}
    \question \textbf{Box and ramp:} Consider a box of mass $m$ being pushed down a ramp, where the pushing force $F$ is always applied parallel to the ground. The vertical distance down the ramp is $y$ and the horizontal distance down the ramp is $x$.
    \begin{figure}[h!]
        
        \centering
        \includegraphics[height=1in]{Images/SimpleRamp.pdf}
    \end{figure}
    \begin{parts}
        \part Draw a free body diagram for the box that has the following four forces: pushing force (F), gravity (G), normal force (N), and friction force (f).
        \part Write a vector $\Delta \vec{r}$ that describes its displacement when it reaches the bottom of the ramp.
        \part For each force on the box, determine if it does positive, negative, or zero work while the box goes down the ramp.
        \part Write an expression for the velocity of the box at the bottom of the ramp in terms of given quantities and the magnitude of the friction force $f$ if it starts sliding from rest.
    \end{parts}
\end{questions}

\section{Potential energy}

\begin{questions}
    \question \textbf{Impact speed}: Consider that you are standing on top of a building  10 m tall, and your friend stands on the ground below. You are both holding identical pumpkins of mass 5 kg.
    \begin{parts}
        \part How much more gravitational energy does your pumpkin have compared to your friend's? Your answer should be in Joules.
        \part How much kinetic energy would the pumpkin have when it reaches the ground if you dropped it from rest?
        \part How much kinetic energy would the pumpkin have when it reaches the ground if you threw it with initial velocity 2 m/s (straight up or down)?
        \part What is the pumpkin's impact speed in these two cases?
        \part Consider air resistance: if you were to go measure this somehow, would you expect your answers in part d to be over or under estimated? Explain why, your reasoning should consider the work done by air resistance.
    \end{parts}
\end{questions}

\section{Springs, universal gravitation, and escape velocity}

% No old question covered spring potential energy or escape velocity.

\end{document}
